\documentclass[11pt]{article}
\usepackage[margin=1in]{geometry}
\usepackage{amsmath,amssymb,amsthm}
\usepackage[T1]{fontenc}
\usepackage{lmodern}
\usepackage{microtype}
\usepackage[colorlinks=true,allcolors=blue]{hyperref}
\newtheorem{theorem}{Theorem}[section]
\newtheorem{lemma}[theorem]{Lemma}
\newtheorem{corollary}[theorem]{Corollary}
\newtheorem{proposition}[theorem]{Proposition}
\title{Sparse pancyclic subgraphs at the full threshold}
\author{Research draft}
\date{10 November 2026}
\begin{document}
\maketitle
\begin{abstract}
Let $\operatorname{Pex}(G)$ be the minimum number of edges above $|V(G)|$ in a spanning pancyclic subgraph of $G$.
We prove that $\operatorname{Pex}(G(n,p))=(1+o(1))\log_2n$ with high probability whenever
$np-\ln n-\ln\ln n\to\infty$.
Writing $w=\min\{np-\ln n-\ln\ln n,\ln\ln n\}$, the construction gives the quantitative upper bound
$\log_2n+(4+o(1))\ln n/\ln w$.
The proof adapts the binary and mixed-radix architecture of Alon and Krivelevich.
Adaptive choices of attachment edges replace endpoint union bounds that require a larger probability gap.
A separate chain of short cycles preserves the binary segment used for long cycle lengths.
\end{abstract}

\paragraph{Scope and attribution.}
Alon and Krivelevich proved the same leading excess bound above an explicit probability
$p_*=(1+O(1/\ln\ln n))\ln n/n$.
We retain their binary encoding, two-stage interval coverage, and independent Hamilton completion.
The changes here are a uniform disjoint-cycle list, equal-depth trees with a fixed linear number of leaves,
and adaptive endpoint selection when the auxiliary mean degree tends to infinity arbitrarily slowly.
The separate short-cycle chain keeps all required short lengths while allowing the auxiliary exposure budget to fit the full threshold.
The external inputs are the rooted tree theorem of Haxell and its stated corollary,
and the Hamilton-path endpoint form of the classical Hamiltonicity threshold.
Precise statements and sources appear below.

This draft was developed in the present research session and underwent three independent internal audits.
The audits checked probability conditioning, rollback after failed attachments, path lengths, Hamilton integration, and cycle-rank accounting.
This is a mathematical research proof, not a Lean formalisation or a publication-priority determination.

\section{Inputs for the threshold construction}
\label{sec:rhost-threshold-tools}

All graphs are simple.
A path's length is its number of edges.
The neighbourhood $N_J(A)$ is external: it contains vertices outside $A$ adjacent to $A$.
The cycle rank of a connected graph is $|E|-|V|+1$.
We write $\ln$ for the natural logarithm and $\log_2$ for the binary logarithm.
An event holds with high probability if its probability tends to one as $n\to\infty$.

\begin{lemma}[Uniform robust expansion]\label{lem:rhost-threshold-expansion}
Suppose $D\to\infty$, $D\le(\ln n)^B$ for a fixed $B$, and $\beta=20\ln D/D$.
Let $R\subseteq[n]$ have size $n-o(n)$ and expose $J\sim G(R,D/n)$.
For any prescribed $r_n=o(n)$, with high probability every $U\subseteq R$ of size at least $n-r_n$
contains a set $U'$ with
\[
 |U'|>|U|-\beta n,\qquad
 |N_{J[U']}(A)|\ge\frac{|A|}{3\beta}
 \quad(0<|A|\le\beta n).
\]
The assertion also holds conditionally on a random $R$ chosen independently of $J$.
\end{lemma}

\begin{proof}
Put $s=\lceil\beta n\rceil$.
The probability that two disjoint $s$-sets in $R$ have no joining edge is at most
\[
 \binom ns^2e^{-Ds^2/n}
 \le\exp\{s[2\ln(en/s)-Ds/n]\}=o(1).
\]
Indeed, $Ds/n\sim20\ln D$ whereas $\ln(en/s)\sim\ln D$.
Assume no such pair exists.
Starting from $U$, repeatedly remove a nonempty set of at most $\beta n$ vertices that violates the desired expansion.
If the total removed set $B$ first reaches size $\beta n$, then $\beta n\le|B|<2\beta n$.
Its external neighbourhood in the original $J[U]$ has size less than $|B|/(3\beta)<2n/3$.
Thus $U\setminus(B\cup N_{J[U]}(B))$ has at least $n/3-r_n-2\beta n$ vertices.
For large $n$, this exceeds $s$.
Taking $s$ vertices there and $s$ vertices in $B$ gives the forbidden pair.
Hence the deletion process removes fewer than $\beta n$ vertices and proves the assertion.
The probability estimate is uniform in $R$.
\end{proof}

The rooted tree theorem below is the form of Haxell's embedding result used in the construction.
It is stated as Theorem~2.2 by Alon and Krivelevich,
\href{https://doi.org/10.1137/23M1598969}{\emph{Sparse Pancyclic Subgraphs of Random Graphs}} (2025).
That statement follows from Haxell, \emph{Tree embeddings}, Journal of Graph Theory 36 (2001), 121--130,
and Balogh, Csaba, Pei, and Samotij,
\href{https://doi.org/10.37236/278}{\emph{Large bounded degree trees in expanding graphs}} (2010).

\begin{theorem}[Rooted tree embedding, external]\label{thm:rhost-rooted-tree}
Let $N,\Delta,k$ be positive integers.
Suppose a graph $J$ satisfies
\[
 |N_J(A)|\ge\Delta|A|+1\quad(0<|A|\le k),
\]
\[
 |N_J(A)|\ge\Delta|A|+N\quad(k<|A|\le2k).
\]
Then every tree with at most $N$ vertices and maximum degree at most $\Delta$
embeds in $J$ with any specified root mapped to any specified vertex of $J$.
\end{theorem}

\begin{lemma}[Trees with many leaves at one depth]\label{lem:rhost-threshold-trees}
Use $\beta$ from Lemma~\ref{lem:rhost-threshold-expansion} and set
\[
 a=\lfloor(5\beta)^{-1}\rfloor,\quad M=\lfloor n/1000\rfloor,
 \quad d=\lceil\log_a M\rceil.
\]
Every expander $J[U']$ supplied by that lemma contains, rooted at any prescribed vertex,
a tree with exactly $M$ leaves at depth $d$, fewer than $3M$ vertices, and maximum degree at most $a+1$.
It also contains a tree formed by joining the roots of two such trees with any positive-length $o(n)$ stem.
\end{lemma}

\begin{proof}
Take $k=\lfloor\beta n/2\rfloor$, $\Delta=\lfloor1/(4\beta)\rfloor$, and $N=\lfloor n/50\rfloor$.
For $0<|A|\le k$, the expansion surplus over $\Delta|A|$ is at least $|A|/(12\beta)>1$.
For $k<|A|\le2k$, it exceeds $n/24>N$.
Thus Theorem~\ref{thm:rhost-rooted-tree} applies.
For sufficiently large $n$, we have $a\ge3$ and $a+1\le\Delta$.

Build complete $a$-ary levels through depth $d-1$.
Give the $a^{d-1}$ vertices at that depth between one and $a$ children each, with exactly $M$ children in total.
This is possible because $a^{d-1}\le M\le a^d$.
All leaves have depth $d$.
The earlier levels have fewer than $2M$ vertices, so the whole tree has fewer than $3M$ vertices.
Two such trees, joined at their roots by a positive-length $o(n)$ stem, have fewer than $6M+o(n)<N$ vertices.
Their maximum degree is at most $a+1$.
Apply the rooted embedding theorem to each required tree.
\end{proof}

\begin{lemma}[A short list of disjoint cycles]\label{lem:rhost-cycle-list}
Fix constants $A,B>0$.
Suppose $D\to\infty$ and $D\le(\ln n)^B$.
In $G(n,D/n)$, with high probability, one can select vertex-disjoint cycles from any fixed list of $O(\ln n)$ lengths,
all between $3$ and $A\ln n$, provided each length has bounded multiplicity.
\end{lemma}

\begin{proof}
Let $X_\ell$ count cycles of length $\ell=O(\ln n)$ and let $Y_\ell$ count ordered intersecting pairs, including identical copies.
Write $\mu_\ell=\mathbb EX_\ell$.
We first establish
\[
 \operatorname{Var}X_\ell,\ \mathbb EY_\ell
 \le C\mu_\ell^2\left(\frac{\ell^2}n+\ell D^{-\ell}\right).
\]
Every proper overlap is a union of paths and isolated vertices.
An overlap with $u$ vertices and $f$ edges has normalized weight $(1+O(\ell^2/n))n^{-u}(n/D)^f$.
An isolated vertex has at most $\ell^2$ descriptions.
A path with $j\ge1$ edges has at most $4\ell^2$ descriptions, each of weight $n^{-1}D^{-j}$.
For several components, division by their order factorial bounds the total by the exponential of the one-component sum.
Thus all proper nonempty overlaps contribute $O(\ell^2/n)$.
The full overlap contributes $(1+O(\ell^2/n))2\ell D^{-\ell}$.
Vertex-disjoint pairs contribute at most $\mu_\ell^2$, proving the variance bound as well.

Put $q=\lfloor\ln n/(10\ln D)\rfloor$.
Chebyshev's inequality, summed over $3\le\ell\le q$, shows that every required short length has any fixed number of copies.
Indeed, $\sum_{\ell\ge3}\ell D^{-\ell}=O(D^{-3})$, and $\mu_\ell\to\infty$ uniformly.
With high probability, all distinct cycles of lengths at most $q$ are vertex-disjoint.
The same proper-overlap enumeration for two possibly different lengths bounds the expected number of intersecting pairs by
\[
 n^{-1}(\ln n)^{O(1)}D^{2q}=o(1).
\]
Choose all prescribed cycles in this range first.

For $q<\ell\le A\ln n$, put $\varepsilon_\ell=C(\ell^2/n+\ell D^{-\ell})\le n^{-1/10+o(1)}$.
Let $M_0=\lceil(\ln n)^3\rceil$.
Chebyshev's inequality gives $X_\ell\ge\mu_\ell/2$ except with probability $O(\varepsilon_\ell)$.
Markov's inequality gives $Y_\ell\le\mu_\ell^2/(8M_0)$ except with probability $O(M_0\varepsilon_\ell)$.
In the graph whose vertices are copies and whose edges join intersecting copies, the elementary independent-set bound is $X_\ell^2/Y_\ell$.
Thus at least $2M_0$ copies are pairwise vertex-disjoint.
A union bound gives this simultaneously for all required lengths.
The prescribed list uses only $O((\ln n)^2)$ vertices in total.
Greedily choose the remaining cycles: each previously used vertex excludes at most one member of any disjoint-copy family.
\end{proof}

The Hamilton-path input below is Lemma~2.1 in the cited Alon--Krivelevich paper.
It follows from the Hamiltonicity threshold and rotation-extension arguments; see Section~4 of
Krivelevich, \href{https://arxiv.org/abs/1507.00205}{\emph{Long paths and Hamiltonicity in random graphs}} (2016).

\begin{lemma}[Hamilton completion, external input]\label{lem:rhost-threshold-hamilton}
If $Nq-\ln N-\ln\ln N\to\infty$, then $G(N,q)$ has, with high probability,
a set $S$ of at least $N/4$ vertices such that every $s\in S$ has a set $T_s$ of at least $N/4$ vertices,
and each $t\in T_s$ is joined to $s$ by a Hamilton path.
Consequently, two external vertices with fresh independent $q$-edges to this graph can be joined through all $N$ vertices with high probability.
\end{lemma}

\begin{proof}[Proof of the consequence]
Expose the induced graph first.
The first external vertex has a neighbour $s\in S$ except with probability at most $e^{-qN/4}$.
Condition on that neighbour.
The second external vertex has a neighbour in $T_s$ with the same failure bound.
These two edges and the stated Hamilton path give the required path.
\end{proof}

\begin{lemma}[Long interval chords]\label{lem:rhost-threshold-windows}
Let a Hamilton cycle on $n$ vertices contain a distinguished arc of length
$2s-1+\ell-1$, where $s\gg n^{0.9}$ and $\ell=o(n^{0.9})$.
Suppose replacing paths inside this arc permits every length reduction from $0$ through $2s-1$.
Let $m=\lfloor n/s\rfloor\ge3$.
In an independent $G(n,D/n)$ with $D\to\infty$, with high probability $m-2$ additional edges
supply cycle intervals that, together with
\[
 [\ell,\ell+2s-1]\quad\text{and}\quad[n-2s+1,n],
\]
cover every integer from $\ell$ through $n$.
\end{lemma}

\begin{proof}
Write $W=n^{0.9}$ and let the distinguished arc have length $b_0=2s+\ell-2$.
For each $3\le i\le m$, there are at least $W^2/100$ vertex pairs whose joining arc contains the distinguished arc
and has length $a_i\in[is-W,is+W]$.
To see this, extend the two ends of the distinguished arc by distances $u,v\ge0$.
Put $t_i=\min\{is-b_0,n-b_0-W/2\}$.
Vary each distance independently in a window of width $W/4$ centered at $t_i/2$.
The lower bound $is-b_0\ge s-o(s)\gg W$ keeps both distances positive.
The upper adjustment leaves a complementary arc of length at least $W/4$.
All these pairs are distinct non-cycle edges and satisfy the required length window.

The probability that some window has no random edge is at most
\[
 m\exp(-Dn^{0.8}/100)=o(1).
\]
Choose one edge from each window.
It closes a cycle on the arc of length $a_i$, and the allowed reductions give
$[a_i+1-(2s-1),a_i+1]$.
Consecutive intervals overlap because $s\gg W$.
The first overlaps $[\ell,\ell+2s-1]$, and the last overlaps $[n-2s+1,n]$ because $ms>n-s$.
This proves the assertion.
\end{proof}

Finally, every connected pancyclic graph with $n+k$ edges has cycle rank $k+1$.
Its distinct cycles give distinct nonzero vectors in its binary cycle space.
Hence $n-2\le2^{k+1}-1$, so $k\ge\log_2(n-1)-1$.

% Independently audited by root, four_regular, and path_stability.
\section{Adaptive attachments at the full pancyclicity threshold}
\label{sec:rhost-full-threshold}

The endpoint failure probability need not be small enough for a direct union bound.
We instead select the attachment edges after exposing each connector tree.
Long base cycles supply spare candidate edges.
For the short cycles, the number of candidates grows with their length.

\begin{theorem}[The leading term at the full threshold]\label{thm:rhost-full-threshold}
Suppose
\[
 np-\ln n-\ln\ln n\longrightarrow\infty.
\]
Then with high probability
\[
 \operatorname{Pex}(G(n,p))=(1+o(1))\log_2 n.
\]
\end{theorem}

\begin{proof}
Put
\[
 w=\min\{np-\ln n-\ln\ln n,\ln\ln n\},\qquad D=w/100000.
\]
Expose five independent graphs with
\[
 p_1=p_2=p_3=p_5=D/n,\qquad
 p_4=(\ln n+\ln\ln n+w/2)/n.
\]
Their union is dominated by $G(n,p)$.
We have $D\to\infty$, without any specified rate.

Use
\[
 \beta=20\ln D/D,\quad a=\lfloor(5\beta)^{-1}\rfloor,\quad
 M=\lfloor n/1000\rfloor,\quad d=\lceil\log_a M\rceil.
\]
Thus $d\sim\ln n/\ln D=o(\ln n)$.
Lemmas~\ref{lem:rhost-threshold-expansion} and~\ref{lem:rhost-threshold-trees} supply the robust expanders and equal-depth leaf trees.
In particular, every permitted reservoir contains an induced expander of size $n-o(n)$.
It contains the required rooted trees and two-tree structures.
Each leaf tree has exactly $M$ leaves at depth $d$ and fewer than $3M$ vertices.
Write $H=\lfloor n/4000\rfloor$ and $\delta=\exp(-D/5000)=o(1)$.
A fresh endpoint has a neighbour in any specified $H$-vertex subset with failure at most $\delta$.
The same bound holds for a specified leaf set after at most $2H$ earlier endpoint queries.

Set
\[
 K_0=\lceil\log_2(4d)\rceil,\quad
 \ell_*=(K_0+2)(d+3),\quad b=\lceil\sqrt{2\ell_*}\rceil,
\]
\[
 c=8b+5,\qquad T=c+4d+3=4d+8b+8,
\]
\[
 K=\left\lfloor\log_2\frac{nw}{100\ln n}\right\rfloor,
 \qquad L_0=2^{K+1}-1,\qquad m=\lfloor n2^{-K}\rfloor.
\]
These parameters are fixed before exposure.
They satisfy
\[
 T,b,d,m=o(\ln n),\quad \ell_*=o(L_0),\quad L_0\gg n^{0.9},
 \quad 2L_0\le\frac{nw}{25\ln n}.
\]
Also, $\ell_*/K\to\infty$, since $D\le\ln\ln n$ and $\ln d\sim\ln\ln n$.

\emph{Initial cycles.}
In $G_1$, select vertex-disjoint bit cycles $C_i$ of lengths $2^i+2$, for $0\le i\le K_0$.
Select a digit base cycle $C_{\rm short}$ of length $c$.
Select one further cycle $S_j$ of each length $3\le j\le T$.
Lemma~\ref{lem:rhost-cycle-list} applies: all lengths are $O(\ln n)$, and their multiplicities are at most three.
The selected cycles use $O(T^2+d)$ vertices.
Also select an edge $f$ outside their vertices.
Such an edge exists with high probability: $G_1$ has $\Omega(nD)$ edges and maximum degree at most $(\ln n)^2$,
while the selected vertices number only $O((\ln n)^2)$.
Let $H_1$ contain all these cycles and $f$.

\emph{An adaptive chain of cycles.}
We describe a connector procedure in $G_2$.
Expose first only edges with both endpoints outside $H_1$.
The universal robust-expansion event then determines every required reservoir expander.
Edges from vertices of $H_1$ to the reservoir remain independent until queried.

Start at a specified endpoint outside the reservoir.
Query at most $H$ fresh candidate neighbours in the current expander.
If a neighbour $z$ exists, embed a leaf tree rooted at $z$.
This first step fails with probability at most $\delta$.
Suppose the next object is an untouched cycle $Q$ of length $j$.
Take a matching of $\lfloor j/2\rfloor$ candidate edges in $Q$.
For a candidate edge $xy$, first query the edges from $x$ to the previous leaf set.
If one exists, it gives an incoming path of length $d+2$ from the preceding endpoint to $x$.
Temporarily exclude this path's internal vertices from the next reservoir.
Then query at most $H$ fresh edges from $y$ to the resulting expander.
If this succeeds, retain the incoming path, select $xy$, and embed the next leaf tree at the chosen neighbour of $y$.
If either query fails, discard that candidate edge and retain no part of its tentative incoming path.
The previous leaf tree remains available for the next candidate.

Each candidate has two endpoints unused in any previous candidate.
The reservoir was exposed before all these endpoint queries.
Even after conditioning on earlier failures and on the tentative incoming path, the edges from the new endpoint $y$ remain fresh.
Thus each conditional candidate failure probability is at most $2\delta$.
A cycle of length $j$ fails with probability at most $(2\delta)^{\lfloor j/2\rfloor}$.
For any list with bounded multiplicity of cycle lengths, their total failure probability is at most
\[
 O\!\left(\sum_{j\ge3}(2\delta)^{\lfloor j/2\rfloor}\right)=o(1).
\]
At the end, one fixed target endpoint attaches to the final leaf set with failure at most $\delta$.
Only the chosen paths, not the unused tree vertices or failed queries, enter the constructed subgraph.

Apply this procedure first to the bit cycles, starting at one endpoint of $f$ and ending at the other.
The selected edge of $C_i$ is denoted $e_i$.
Together with $f$ and the $K_0+2$ connectors, these edges form a cycle $C_*$ of length $\ell_*$.
Each other arc of $C_i$ replaces $e_i$ with gain exactly $2^i$.

Fix an edge $e_{\rm short}$ of $C_{\rm short}$.
Apply the procedure again to $S_3,\ldots,S_T$.
Start at the endpoint of $f$ used for the first bounded query, and end at one endpoint of $e_{\rm short}$.
The starting endpoint has only $H$ earlier queries, so another $H$ fresh candidate edges remain available.
This second chain is disjoint from the first construction except at its designated endpoint.
It contains all single-cycle blocks needed for the short lengths.

\emph{Adaptive binary and digit detours.}
Expose the $G_3$ reservoir edges outside the entire graph constructed so far.
All its base-cycle endpoint edges remain fresh.
For each $i>K_0$, embed two depth-$d$ leaf trees whose roots are joined by a stem of length
\[
 2^i+1-2d-2=2^i-2d-1>0.
\]
The total tree size is below $n/50$.
Choose an unused candidate edge $xy$ of $C_*$, and query its two endpoints against the two leaf sets.
If both attachments succeed, the resulting detour has length $2^i+1$.
Its replacement gain over $xy$ is $2^i$.
If an attachment fails, abandon that candidate edge and try another with two fresh endpoints.

Preselect $2(K-K_0)$ pairwise vertex-disjoint candidate edges of $C_*$,
avoiding $f$ and all endpoints of $e_0,\ldots,e_{K_0}$.
There are enough because $\ell_*/K\to\infty$.
Every conditional failure probability is at most $2\delta$.
The expected number of failures in this entire candidate pool is at most $4\delta(K-K_0)$.
If the pool is exhausted before completing all detours, at least $K-K_0$ candidates have failed.
Markov's inequality shows that this event has probability at most $4\delta=o(1)$.
After each success, retain only its detour path and re-form the next reservoir.

The digit base cycle has at least $4b$ pairwise vertex-disjoint candidate edges avoiding the endpoints of $e_{\rm short}$.
Use them to obtain $2b$ successful detours, in two groups of $b$.
The desired detour lengths are
\[
 2d+3+j,\qquad 2d+3+jb\qquad(0\le j<b).
\]
The two-tree stems have lengths $1+j$ and $1+jb$, so they are positive.
The same expected-failures argument succeeds with probability $1-o(1)$.
The designated edges need not be consecutive on the base cycle.
Choosing one detour from each group gives every cycle length in
\[
 [T+1,T+b^2].
\]
Indeed, the length is $c-2+(2d+3+j_0)+(2d+3+bj_1)=T+1+j_0+bj_1$.
This interval overlaps $[\ell_*,\ell_*+L_0]$, since $T=o(\ell_*)$ and $b^2\ge2\ell_*$.

\emph{Hamilton completion and final coverage.}
Delete $f$, $e_{\rm short}$, and the selected edge of each $S_j$ when choosing a spanning path through the constructed blocks.
Use every bit detour, every digit detour, and the long arc of each $S_j$ in this path.
The separate short-cycle chain joins one end of the binary path to one end of the digit path.
Thus all constructed vertices lie on one spanning path.
The binary path is a contiguous segment of it.
All deleted edges remain present as chords in the constructed subgraph.

The total used vertex count is at most $2L_0$.
The binary detours contribute $L_0+O(\ell_*)$ vertices.
The short-cycle chain and digit gadget contribute only a polynomial in $\ln n$.
For the remaining vertex set $X$,
\[
 |X|p_4-\ln|X|-\ln\ln|X|
 \ge(1/2-1/25-o(1))w\longrightarrow\infty.
\]
Lemma~\ref{lem:rhost-threshold-hamilton} therefore joins the two remaining endpoints through all vertices of $X$.
This is the completion input from Alon--Krivelevich Lemma~2.1 and Claim~4.7.
Its endpoint failure probability is $o(1)$.
This gives a Hamilton cycle.

The bit cycles give every length in $[\ell_*,\ell_*+L_0]$.
Replacing bit detours on the Hamilton cycle gives every length in $[n-L_0,n]$.
The individual cycles $S_j$ and the digit gadget cover all lengths below the first interval.
Use $G_5$ only for the long interval-covering chords.
Lemma~\ref{lem:rhost-threshold-windows} applies because the binary segment is contiguous, has length $L_0+\ell_*-1$,
and $\ell_*=o(n^{0.9})=o(L_0)$.
Each required window has at least $n^{1.8}/100$ candidate edges.
The failure probability is at most $m\exp(-Dn^{0.8}/100)=o(1)$.
Their overlapping length intervals cover the remaining gap.

The constructed graph has excess at most
\[
 K+2b+T+m+O(1)=\log_2 n+o(\ln n).
\]
The short-cycle chain contributes $T-2$ to cycle rank.
Failed queries contribute no retained edges.
The universal cycle-count lower bound gives the asserted equality to leading order.
\end{proof}

\begin{corollary}[A quantitative full-threshold error]\label{cor:rhost-full-threshold-error}
Write $\omega=np-\ln n-\ln\ln n\to\infty$ and $w=\min\{\omega,\ln\ln n\}$.
Then with high probability
\[
 \operatorname{Pex}(G(n,p))
 \le\log_2n+(4+o(1))\frac{\ln n}{\ln w}.
\]
\end{corollary}

\begin{proof}
The preceding construction has $\ln a=(1+o(1))\ln w$, so $d\sim\ln n/\ln w$.
Moreover, $b/d=O(\sqrt{K_0/d})=o(1)$ and $m/d=O(\ln w/w)=o(1)$.
Therefore $T=4d+o(d)$, and the displayed excess bound gives the result.
\end{proof}


\paragraph{Finite verification.}
The script \texttt{check\_full\_threshold\_integration.py} checks the finite gadget topology,
arbitrary disjoint choices of attachment edges, exact binary and digit length identities,
and the final Hamilton cycle and excess count.
It passes 16 gadget cases, 3072 binary-cycle checks, 1784 digit-cycle checks, and 16 spanning Hamilton-cycle checks.
These checks do not replace the asymptotic probability proofs.

\begin{thebibliography}{9}
\bibitem{AK}
Y. Alon and M. Krivelevich,
\emph{Sparse Pancyclic Subgraphs of Random Graphs},
SIAM Journal on Discrete Mathematics \textbf{39} (2025), 562--574.
\href{https://doi.org/10.1137/23M1598969}{doi:10.1137/23M1598969}.
The authors' checked manuscript is dated 29 November 2024:
\url{https://www.math.tau.ac.il/~krivelev/pancyclic-rg.pdf}.

\bibitem{Haxell}
P. E. Haxell, \emph{Tree embeddings},
Journal of Graph Theory \textbf{36} (2001), 121--130.
\url{https://doi.org/10.1002/1097-0118(200103)36:3\%3C121::AID-JGT1000\%3E3.0.CO;2-U}.

\bibitem{BCPS}
J. Balogh, B. Csaba, M. Pei, and W. Samotij,
\emph{Large bounded degree trees in expanding graphs},
Electronic Journal of Combinatorics \textbf{17} (2010), Research Paper~6.
\href{https://doi.org/10.37236/278}{doi:10.37236/278}.

\bibitem{Krivelevich}
M. Krivelevich, \emph{Long paths and Hamiltonicity in random graphs},
in N. Fountoulakis and D. Hefetz (eds.),
\emph{Random Graphs, Geometry and Asymptotic Structure},
London Mathematical Society Student Texts \textbf{84}, Cambridge University Press, 2016, pp.~4--27.
\url{https://arxiv.org/abs/1507.00205}.
\end{thebibliography}
\end{document}
